\chapter*{Conclusion générale}
\addcontentsline{toc}{chapter}{Conclusion générale}
\markboth{Conclusion générale}{Conclusion générale}

Ce mémoire a présenté la conception et l'implémentation d'un système distribué de mitigation collaborative des attaques \gls{ddos}, reposant sur une coalition pair-à-pair de n\oe{}uds de protection hétérogènes. Face aux limites structurelles des approches centralisées --- points de défaillance uniques, coûts élevés et absence de mutualisation des ressources --- le système proposé offre une alternative décentralisée où chaque participant prend ses décisions de manière autonome tout en contribuant à la défense collective.

Le premier chapitre a posé les fondements conceptuels~: taxonomie des attaques \gls{ddos}, mécanismes d'infrastructure de mitigation existants et principes des architectures pair-à-pair. Le deuxième chapitre a analysé les approches collaboratives existantes (\gls{dots}, DDoS Clearing House, NaWas/NBIP, DefCOM et PeerTrust) et identifié leurs lacunes en matière de mécanismes de confiance et de sélection intelligente des pairs. Le troisième chapitre a présenté la conception du système~: la conception préliminaire a défini les exigences et l'architecture globale de la coalition~; la conception technique détaillée a formalisé le modèle de données relationnel en quatorze tables, le mécanisme de confiance PeerTrust adapté, l'algorithme de sélection \gls{wsm} avec pondération \gls{ahp}, et les protocoles de sécurité mTLS et JWT RS256. Le quatrième chapitre a détaillé l'implémentation et l'évaluation du prototype. Les tests de scalabilité conduits sur une coalition simulée de 20 à 100~n\oe{}uds ont confirmé un taux de succès de 100\,\% pour les opérations d'enregistrement et de heartbeat à tous les paliers, un temps de sélection \gls{wsm} inférieur à 33~ms quasi-constant quelle que soit la taille de la coalition, et des temps de réponse heartbeat restant en dessous de 500~ms jusqu'à 100~n\oe{}uds simultanés.

Les contributions principales de ce travail sont au nombre de trois. La première concerne l'\textbf{adaptation du modèle PeerTrust avec crédibilité historique}~: la crédibilité $Cr_i$ est désormais capturée à la clôture de chaque session d'aide et persistée en base de données, rendant le calcul de $T(u)$ fidèle à la réalité temporelle du réseau. La deuxième contribution porte sur la \textbf{sélection multicritère des pairs par WSM-AHP}~: l'algorithme combine quatre critères (capacité, latence, confiance, réciprocité) avec des poids validés (CR~$= 1{,}6\,\%$), permettant à tous les pairs éligibles de participer simultanément. La troisième contribution réside dans la \textbf{sécurité asymétrique sans secret partagé}~: JWT RS256 combiné à mTLS garantit que la compromission d'un n\oe{}ud n'expose pas les autres membres de la coalition.

Ce travail ouvre plusieurs perspectives d'amélioration. En premier lieu, la mise en place d'un \textbf{mécanisme de découverte automatique} constitue une priorité~: un protocole de bootstrap similaire aux architectures \gls{dht} permettrait à un n\oe{}ud de rejoindre la coalition en ne connaissant qu'un seul pair initial. La \textbf{détection des comportements malveillants avancés} mérite également d'être approfondie, notamment les attaques Sybil et le free-riding coordonné, par l'intégration d'un module de détection d'anomalies comportementales. Sur le plan de la scalabilité, le \textbf{passage à l'échelle au-delà de 100~n\oe{}uds} nécessitera l'ajout d'un cache distribué (Redis) pour les scores de confiance. À plus long terme, une \textbf{passerelle vers \gls{dots}} (RFC~8782) permettrait à la coalition de s'interfacer avec les prestataires de mitigation commerciaux. Enfin, une \textbf{évaluation en environnement réel} sur une infrastructure géographiquement distribuée permettrait de valider les résultats obtenus en simulation et de mesurer l'impact de la latence réseau réelle sur les mécanismes de recalcul de confiance.

En définitive, ce travail démontre la faisabilité d'une défense collaborative distribuée contre les attaques \gls{ddos}, reposant sur des mécanismes de confiance formels et une sélection intelligente des pairs. Il constitue une base solide pour le développement d'une infrastructure de coalition opérationnelle, adaptée aux contraintes hétérogènes des organisations participantes.
